\documentclass[a4paper]{article}
\usepackage{amsfonts}
\usepackage{amsmath}
\usepackage{amssymb}
\usepackage{graphicx}     

\begin{document}
  
\noindent \large Auteur: CouldBeMathijs \\
\noindent \large Studierichting: Informatica\\
\noindent \large Oefeningenreeks 1D nr. 20.5\\

\medskip

\normalsize

\begin{eqnarray*}
z^4 + 4 &=& 0 \\
z^4 &=& -4 \\
-4 &=& 4\left(\cos(\pi + 2k\pi) + i\sin(\pi + 2k\pi)\right)\\
z_k &=& \sqrt[4]{4}
\left(\cos\left(\dfrac{\pi + 2k\pi}{4}\right) + i\sin\left(\dfrac{\pi + 2k\pi}{4}\right)\right) \\
&=& \sqrt{2}\left(\cos\left(\dfrac{\pi}{4} + \dfrac{k\pi}{2}\right) + i\sin\left(\dfrac{\pi}{4} + \dfrac{k\pi}{2}\right) \right), \text{met } k=0,1,2,3
\end{eqnarray*}

Complexe nulpunten:

\begin{eqnarray*}
z_0 &=& 1 + i \\
z_1 &=& -1 + i \\
z_2 &=& -1 - i \\
z_3 &=& 1 - i
\end{eqnarray*}

Ontbinding van de veelterm in $\mathbb{R}$:\\

\begin{eqnarray*}
    z^4 + 4 &=& ((z-z_0) (z-z_1)) ((z-z_2) (z-z_3)\\
    &=& (z^2 - 2z + 2)(z^2 + 2z + 2) \text{ waar elke factor een discriminant kleiner dan 0 heeft.}
\end{eqnarray*}




\begin{thebibliography}{999}
%\bibitem{a} Referentie
\end{thebibliography}
\end{document} 
